\ssscp{Distribution of reflexes of PMP *(h)aŋəluq ‘wind’}{15-04-02-02}
Reflexes of PMP *haŋin ‘wind’ are found throughout the Malayo-Polynesian world.
\citet[110]{Blust2008} also reconstructs Proto Sumba-Havu **ŋəlu
‘wind’ based on Kambera \it{ŋilu} ‘wind’
and Havu \it{ŋəlu} ‘wind’.
While that is entirely reasonable,
it also does not show the fuller picture of related terms straddling the Blust line.
The data in \trf{15-02} are from \citet{Zobel2020},
supplemented from our own sources.

\begin{table}[p]
	\centering\tcp{Distribution of reflexes of *(h)aŋəluq ‘wind’}{15-02}
	\st{2pt}\begin{tabular}{llll}\lsptoprule
Language	&	Form	&	Node	&	Source	\\
PMP	&	*(h)aŋəluq	&	PMP	&	\cite{Zobel2020}	\\	\midrule
Chamorro	&	{\hr}mɑŋloʔ	&	\tsc{Chamorro}	&	\cite{Topping1975}	\\	\midrule
Maranao	&	{\hr}ndoʔ	&	\tsc{Philippine}	&	\cite{McKaughanMacaraya1967}	\\	\midrule
Lemolang	&	{\hr}oŋalu	&	\tsc{S. Sulawesi}	&	\cite[115]{GrimesGrimes1987}	\\	\midrule
P.N. Kaili-Wolio	&	**ŋoluʔ	&	\tsc{Celebic}	&	\cite{Zobel2020}	\\
Uma	&	{\hr\hr}ŋoluʔ	&	\tsc{C. Sulawesi}	&	\cite{Zobel2020}	\\
Kulawi	&	{\hr\hr}ŋolu	&	\tsc{C. Sulawesi}	&	\cite{Zobel2020}	\\
Pamona	&	{\hr\hr}ŋoyu	&	\tsc{C. Sulawesi}	&	\cite[115]{GrimesGrimes1987}	\\
P.S. Kaili-Wolio	&	**ŋalu	&	\tsc{Celebic}	&	\cite{Zobel2020}	\\
Wolio	&	{\hr\hr}ŋalu	&	\tsc{Wotu-Wolio}	&	\cite{Zobel2020}	\\
Loa{\Q}	&	{\hr\hr}ŋalu	&	\tsc{Wotu-Wolio}	&	\cite{Zobel2020}	\\
Laiyolo	&	{\hr\hr}ŋalu	&	\tsc{Wotu-Wolio}	&	\cite[115]{GrimesGrimes1987}	\\
Taloki	&	{\hr\hr}ŋalu	&	\tsc{Bungku-Tolaki}	&	\cite[136]{Mead1999}	\\
Kalisusu	&	{\hr\hr}ŋalu	&	\tsc{Bungku-Tolaki}	&	\cite[136]{Mead1999}	\\
Koroni	&	{\hr\hr}ŋalu	&	\tsc{Bungku-Tolaki}	&	\cite[136]{Mead1999}	\\	\midrule
Kambera	&	{\hr\hr}ŋilu	&	\tsc{Bima-Lembata}	&	\cite[348]{Onvlee1984}	\\
Anakalang	&	{\hr\hr}ŋilu	&	\tsc{Bima-Lembata}	&	\cite[348]{Onvlee1984}	\\
Mamboru	&	{\hr\hr}ŋɐlu	&	\tsc{Bima-Lembata}	&	\cite[348]{Onvlee1984}	\\
Havu	&	{\hr\hr}ŋəlu	&	\tsc{Bima-Lembata}	&	\cite{Grimes2008}	\\
Dhao	&	{\hr\hr}ŋəlu	&	\tsc{Bima-Lembata}	&	\cite{GrimesRanohAplugi2008}	\\	\midrule
Welaun	&	{\hr\hr}kolu	&	\tsc{Central Timor}	&	\cite[43]{Edwards2019}	\\
Tokodede	&	{\hr\hr}gelu	&	\tsc{Central Timor}	&	\cite[53]{Edwards2019}	\\
NW. Mambae	&	{\hr\hr}gelo	&	\tsc{Central Timor}	&	\cite[231]{Fogaca2017}	\\
C. Mambae	&	{\hr\hr}kiul-a	&	\tsc{Central Timor}	&	\cite[231]{Fogaca2017}	\\
S. Mambae	&	{\hr\hr}kelu {\tl} keul	&	\tsc{Central Timor}	&	\cite{GrimesMarcalFereira2014}	\\	\midrule
Selaru	&	{\hr\hr}esw	&	\tsc{Timor-Babar}	&	\cite{CowardND}; *ŋ > {\0}, *l > \it{s}	\\
Seluwasan	&	{\hr\hr}nedw-e	&	\tsc{Timor-Babar}	&	\cite{Conner2012}; *l > \it{d}	\\
Makatian	&	{\hr\hr}neʤw-e	&	\tsc{Timor-Babar}	&	\cite{Conner2012}; *l > \it{ʤ}	\\		\lspbottomrule
	\end{tabular}
\end{table}

%dew
%embun
%BAL ondu -na
%DAM ngolu +DMO
%TIT
%ondu
%PEN ngolu -nyo
%TAl
%ngolu
%AMP ngolu
%LAU ngolu -nye
%TIA
%ngolu
%DON ngolu
%TOL ngolu
%BOA ngou'
%Himmelmann 2001 pdf.p. 124

First note that reflexes of this term are found in widely dispersed
subgroups on both sides of the Blust line.
In some cases **ŋəlu
and related terms appear to be replacement innovations within
whole subgroups, and in others they co-exist side-by-side with reflexes of PMP
*haŋin ‘wind’ in sporadic distribution.
For instance, Makatian and Seluwasan (\tsc{Timor-Babar})
have reflexes of **ŋəlu meaning `wind'
alongside *haŋin > \it{anin} `year'.
Similarly, while most \tsc{Central Timor} languages
have reflexes of **ŋəlu for `wind',
Kemak has *haŋin > \it{agi} ‘wind’.\footnote{
		In some languages of northern Sulawesi 
		cognates of **ŋəlu mean `dew'; e.g.
		Dampelas, Totoli \it{ŋolu} `dew', Boano
		\it{ŋouʔ} `dew' \citep[124f]{Himmelmann2001}.}

If we knew of the reflexes of these terms only in Wallacean subgroups,
we might be tempted to reconstruct a term for PCMP.
But its systematic presence in \tsc{Celebic} languages
and further afield prevents that.
Its (apparent) absence from other Wallacean subgroups also
mitigates against doing so.
A map showing the distribution of cognates of *(h)aŋəluq in
Wallacea and neighbouring regions is given in \frf{15-wind}.

\begin{figure}[h]\centering
	\includegraphics[width=\columnwidth]{wind.pdf}
	\fcp{Distribution of *(h)aŋəluq `wind'}{15-wind}
\end{figure}